\chapter{运算放大器}\label{ch:b1:operational-amplifier}

粘在桥梁大梁上的应变片，卡车经过时其电阻变化千分之一；
而它所在的\omterm{prop:b1:dc-circuits:wheatstone}{惠斯通电桥}
（\cref{ch:b1:dc-circuits}）把这一变化转换为两毫伏的\omterm{def:b1:dc-circuits:voltage}{电压}。
两毫伏的\omterm{def:b1:wave-propagation:signal}{信号}无法驱动任何东西。在应变片与
那套将桥梁封闭以禁止通行的报警装置之间，
有一块指甲盖大小的芯片，售价只需几分钱，它能把\omterm{def:b1:wave-propagation:signal}{信号}放大一千倍、
对其进行滤波、与阈值比较，并点亮一盏灯——四项任务，一个元件，
四种不同的接线方式。本章通过理想模型引入
\omterm{def:b1:operational-amplifier:opamp}{运算放大器}，推导每台仪器都会用到的几种基本电路，
并说明
移除反馈后会发生什么。

\section{元件及其理想模型}

\begin{definition}[运算放大器]\label{def:b1:operational-amplifier:opamp}
\index{运算放大器}\index{饱和电压}\index{开环增益}
\emph{运算放大器}是一种集成电路，有两个输入端——
电势为 $V_+$ 的\emph{同相}输入端与
电势为 $V_-$ 的\emph{反相}输入端——以及电势为 $V_s$ 的一个输出端，
由两路电源 $\pm V_{\mathrm{cc}}$ 供电（通常为
$\pm\qty{15}{V}$），图中一般不画出这两路电源。它放大
两输入端之差 $\epsilon = V_+ - V_-$：
\[
V_s = \mu\,\epsilon \quad\text{只要 } \abs{V_s} < V_{\mathrm{sat}},
\]
其中\emph{开环增益} $\mu$ 的量级为 $10^5$，
\emph{饱和\omterm{def:b1:dc-circuits:voltage}{电压}} $V_{\mathrm{sat}}$ 略低于
$V_{\mathrm{cc}}$。当 $\mu\abs\epsilon$ 将超过 $V_{\mathrm{sat}}$ 时，
输出被钳位在 $+V_{\mathrm{sat}}$ 或 $-V_{\mathrm{sat}}$，
此状态称为\emph{饱和}状态。
\end{definition}

\begin{definition}[理想运算放大器]\label{def:b1:operational-amplifier:ideal}
\index{理想运算放大器}\index{虚短路}\index{线性区}
\emph{理想}模型假设：（i）两输入端均无电流流入，
$i_+ = i_- = 0$（输入\omterm{def:b1:sinusoidal-impedance:impedance}{阻抗}无穷大）；（ii）输出端的\omterm{def:b1:dc-circuits:voltage}{电压}
不受所流出电流的影响（输出\omterm{def:b1:sinusoidal-impedance:impedance}{阻抗}为零）；（iii）
增益无穷大，$\mu \to \infty$。在\emph{线性区}（$\abs{V_s}
< V_{\mathrm{sat}}$）中，第三条假设迫使
\[
\epsilon = V_+ - V_- = 0 :
\]
两输入端电势相等（称为“虚短路”），
但两者之间并无电流通过。
\end{definition}

\begin{omfigure}
\begin{tikzpicture}[font=\small]
  \node[op amp] (oa) at (0,0) {};
  \draw (oa.-) -- ++(-0.8,0) node[left] {$V_-$};
  \draw (oa.+) -- ++(-0.8,0) node[left] {$V_+$};
  \draw (oa.out) -- ++(0.6,0) node[right] {$V_s$};
  \draw (oa.up) -- ++(0,0.3) node[above, font=\footnotesize] {$+V_{\mathrm{cc}}$};
  \draw (oa.down) -- ++(0,-0.3) node[below, font=\footnotesize] {$-V_{\mathrm{cc}}$};
  \begin{scope}[xshift=6.2cm]
  \begin{axis}[omaxis, axis y line=left, width=6.4cm, height=4.6cm, xlabel={$\epsilon$}, ylabel={$V_s$},
      xmin=-2.2, xmax=2.2, ymin=-1.5, ymax=1.5, xtick={-0.15,0.15}, xticklabels={,},
      ytick={-1,1}, yticklabels={$-V_{\mathrm{sat}}$,$+V_{\mathrm{sat}}$}, anchor=origin, at={(0,0)}]
    \addplot[omThm, thick, domain=-2.2:-0.15, samples=2] {-1};
    \addplot[omThm, thick, domain=-0.15:0.15, samples=2] {x/0.15};
    \addplot[omThm, thick, domain=0.15:2.2, samples=2] {1};
    \node[omThm, font=\footnotesize] at (axis cs:1.3,0.6) {饱和};
    \node[omThm, font=\footnotesize] at (axis cs:-1.2,-0.6) {饱和};
    \node[omThm, font=\footnotesize, right] at (axis cs:0.2,-0.55) {线性区，斜率 $\mu$};
  \end{axis}
  \end{scope}
\end{tikzpicture}

{\small \omterm{def:b1:operational-amplifier:opamp}{运算放大器}的符号（带两个输入端），以及其
传输特性 $V_s(\epsilon)$——一段陡峭的线性段，
斜率为 $\mu \sim 10^5$（图中画得远比实际平缓），随后
在 $\pm V_{\mathrm{sat}}$ 处进入饱和。在理想模型中，线性段
是竖直的：$\epsilon = 0$。}
\end{omfigure}

\begin{proposition}[负反馈与线性区]\label{prop:b1:operational-amplifier:feedback}
\index{负反馈}\index{正反馈}
若输出端经由某一网络接回\emph{反相}输入端（\omterm{prop:b1:operational-amplifier:feedback}{负反馈}），
则电路有一个稳定的线性\omterm{met:b1:dc-circuits:loadline}{工作点}，
位于 $\epsilon = 0$ 处，且适用理想放大器的规则。若
输出端接回同相输入端（\omterm{prop:b1:operational-amplifier:feedback}{正反馈}），或根本不接反馈，则放大器饱和：
$V_s = \pm V_{\mathrm{sat}}$，
取决于 $\epsilon$ 的符号。
\end{proposition}

\begin{proof}[部分证明]
在\omterm{prop:b1:operational-amplifier:feedback}{负反馈}下，设 $V_s$ 略高于其平衡值；经由反馈网络，
$V_-$ 升高，$\epsilon$ 减小，放大器将 $V_s$ 拉回，
扰动被修正。在\omterm{prop:b1:operational-amplifier:feedback}{正反馈}下，同样的扰动抬高 $V_+$，
使 $\epsilon$ 增大，
并推动 $V_s$ 进一步远离，直至饱和。
完整论证（包括放大器的一阶模型 $\mu(j\omega)
= \mu_0/(1 + j\omega/\omega_0)$，以及由此得到的微分方程的稳定性）
属于本科二年级分册的内容，
此处只需使用上述规则。
\end{proof}

\begin{remark}[实际放大器]\label{rem:b1:operational-amplifier:real}
\index{增益带宽积}\index{压摆率}\index{输入失调电压}
实际中有三点偏离理想模型的情况需要注意。\omterm{def:b1:operational-amplifier:opamp}{开环增益}随频率下降，
关系为 $\mu_0/(1 + jf/f_0)$，其中 $f_0 \sim \qty{10}{Hz}$：
乘积 $\mu_0 f_0 = f_T \sim \qty{1}{MHz}$（即\emph{\omterm{rem:b1:operational-amplifier:real}{增益带宽
积}}）即为跟随器可用的带宽，
闭环增益为 $G$ 的电路其带宽为 $f_T/G$。输出\omterm{def:b1:dc-circuits:voltage}{电压}的变化速度
不能超过\emph{\omterm{rem:b1:operational-amplifier:real}{压摆率}}，约为 $\sim\qty{1}{V/\micro s}$。另外，
一个较小的\emph{\omterm{rem:b1:operational-amplifier:real}{输入失调电压}}（$\sim\qty{1}{mV}$）会像\omterm{def:b1:wave-propagation:signal}{信号}一样被放大，
对于毫伏量级的输入而言，若不加校正，这将是致命的。
\end{remark}

\section{基本线性电路}

\begin{method}[线性区中理想运算放大器电路的分析]\label{met:b1:operational-amplifier:method}
\begin{enumerate}
  \item 检查反馈是否接到反相输入端。
  \item 令 $i_+ = i_- = 0$：输入节点服从基尔霍夫节点定律，
        即流入放大器的电流为零（分压与米尔曼
        定理均适用）。
  \item 令 $V_+ = V_-$，解出 $V_s$。
  \item 事后检验 $\abs{V_s} < V_{\mathrm{sat}}$ 是否成立，若不成立，
        则放大器已饱和，结果为 $\pm V_{\mathrm{sat}}$。
\end{enumerate}
\end{method}

\begin{proposition}[跟随器、同相放大器与反相放大器]\label{prop:b1:operational-amplifier:basic}
\index{电压跟随器}\index{同相放大器}\index{反相放大器}\index{虚地}
\begin{itemize}
  \item \emph{跟随器}：输出端接 $V_-$，输入端接在 $V_+$ 上：
        $V_s = V_e$；输入\omterm{def:b1:sinusoidal-impedance:impedance}{阻抗}无穷大，输出\omterm{def:b1:sinusoidal-impedance:impedance}{阻抗}为零。
  \item \emph{\omterm{prop:b1:operational-amplifier:basic}{同相放大器}}：输出端接 $V_-$，经
        由 $R_2$（反馈）与 $R_1$（接地）构成的\omterm{prop:b1:dc-circuits:dividers}{分压器}，输入接在 $V_+$ 上：
        \[
        V_s = \left(1 + \frac{R_2}{R_1}\right)V_e .
        \]
  \item \emph{\omterm{prop:b1:operational-amplifier:basic}{反相放大器}}：$V_+$ 接地，输入\omterm{def:b1:wave-propagation:signal}{信号}经
        $R_1$ 接到 $V_-$，反馈电阻 $R_2$ 从输出端接回 $V_-$：
        \[
        V_s = -\frac{R_2}{R_1}\,V_e ,
        \]
        此时 $V_-$ 保持地电势（称为\emph{\omterm{prop:b1:operational-amplifier:basic}{虚地}}），
        输入\omterm{def:b1:sinusoidal-impedance:impedance}{阻抗}为 $R_1$。
\end{itemize}
\end{proposition}

\begin{proof}
跟随器：$V_- = V_s$，$V_+ = V_e$；由 $\epsilon = 0$ 得 $V_s = V_e$。
\omterm{prop:b1:operational-amplifier:basic}{同相放大器}：$V_-$ 处无电流流入，故由分压关系得 $V_- =
V_sR_1/(R_1 + R_2)$，令其等于 $V_+ = V_e$。\omterm{prop:b1:operational-amplifier:basic}{反相放大器}：$V_- = V_+ = 0$；
在 $V_-$ 处列节点方程（取 $i_- = 0$）：$(V_e - 0)/R_1 + (V_s - 0)/R_2 = 0$。
\omterm{def:b1:wave-propagation:signal}{信号}源提供的电流为 $V_e/R_1$，输入\omterm{def:b1:sinusoidal-impedance:impedance}{阻抗}为 $R_1$。
\end{proof}

\begin{omfigure}
\begin{tikzpicture}[font=\small]
  % follower (+ input up)
  \node[op amp, noinv input up] (a) at (0,0) {};
  \draw (a.+) -- ++(-0.6,0) node[left] {$V_e$};
  \draw (a.out) -- ++(0.5,0) node[right] {$V_s$};
  \draw (a.out) ++(0.25,0) -- ++(0,-1.3) -| ($(a.-)+(-0.4,0)$) -- (a.-);
  \node[below, font=\footnotesize] at (0,-1.7) {跟随器：$V_s = V_e$};
  % non-inverting (+ input up)
  \begin{scope}[xshift=5.6cm]
  \node[op amp, noinv input up] (b) at (0,0) {};
  \draw (b.+) -- ++(-0.8,0) node[left] {$V_e$};
  \draw (b.out) -- ++(0.6,0) node[right] {$V_s$};
  \coordinate (j) at ($(b.-)+(-0.5,0)$);
  \draw (b.-) -- (j);
  \coordinate (f) at ($(j)+(0,-1.0)$);
  \coordinate (o) at ($(b.out)+(0.3,0)$);
  \draw (o) -- (o |- f) to[R=$R_2$] (f) -- (j);
  \draw (f) to[R=$R_1$] ++(0,-1.6) node[ground]{};
  \node[below, font=\footnotesize] at (0,-3.7) {同相放大器：$V_s = (1 + R_2/R_1)V_e$};
  \end{scope}
  % inverting (- input up)
  \begin{scope}[xshift=11.3cm]
  \node[op amp] (c) at (0,0) {};
  \draw (c.+) -- ++(-0.4,0) node[ground]{};
  \coordinate (m) at ($(c.-)+(-0.5,0)$);
  \draw (c.-) -- (m) to[R=$R_1$] ++(-1.5,0) node[left] {$V_e$};
  \draw (c.out) -- ++(0.6,0) node[right] {$V_s$};
  \coordinate (f2) at ($(m)+(0,1.3)$);
  \coordinate (o2) at ($(c.out)+(0.3,0)$);
  \draw (o2) -- (o2 |- f2) to[R=$R_2$] (f2) -- (m);
  \node[below, font=\footnotesize] at (0,-1.7) {反相放大器：$V_s = -(R_2/R_1)V_e$};
  \end{scope}
\end{tikzpicture}

{\small 三种基本放大器。在每一种中，输出端都接回
反相输入端，放大器工作在\omterm{def:b1:operational-amplifier:ideal}{线性区}，增益
仅由电阻之比决定——而不取决于放大器自身
巨大且不确定的 $\mu$。}
\end{omfigure}

\begin{example}[为何使用跟随器]\label{ex:b1:operational-amplifier:follower}
一个戴维南电阻为 \qty{10}{k\ohm} 的传感器，接上
\qty{1}{k\ohm} 负载后，只能输出其开路\omterm{def:b1:dc-circuits:voltage}{电压}的 $1/11$
（见\cref{rem:b1:dc-circuits:loading}）。接入跟随器后，它几乎不从
传感器汲取电流，并以零输出\omterm{def:b1:sinusoidal-impedance:impedance}{阻抗}驱动负载，
使满\omterm{def:b1:dc-circuits:voltage}{电压}得以传递。跟随器是一种\emph{\omterm{def:b1:sinusoidal-impedance:impedance}{阻抗}
变换器}——它还使得
\cref{prop:b1:filters-transfer-functions:cascade}中的级联规则严格成立。
\end{example}

\begin{proposition}[求和放大器与差分放大器]\label{prop:b1:operational-amplifier:sumdiff}
\index{求和放大器}\index{差分放大器}
输入 $V_1, V_2$ 分别经 $R_1, R_2$ 接到反相节点，
反馈电阻为 $R_f$，$V_+$ 接地：
\[
V_s = -R_f\left(\frac{V_1}{R_1} + \frac{V_2}{R_2}\right)
\quad (\text{求和放大器}).
\]
若 $V_1$ 经 $R_1$ 接到 $V_-$（反馈为 $R_2$），$V_2$ 经
$R_1$ 接到 $V_+$，并经 $R_2$ 将 $V_+$ 接地：
\[
V_s = \frac{R_2}{R_1}\,(V_2 - V_1)
\quad (\text{差分放大器}).
\]
\end{proposition}

\begin{proof}
求和：在\omterm{prop:b1:operational-amplifier:basic}{虚地}处列节点方程，$V_1/R_1 + V_2/R_2 + V_s/R_f =
0$。差分：$V_+ = V_2R_2/(R_1 + R_2)$（分压）；在 $V_-$ 处列节点方程：
$(V_1 - V_-)/R_1 + (V_s - V_-)/R_2 = 0$，于是 $V_s = V_-(1 + R_2/R_1) -
V_1R_2/R_1$；令 $V_- = V_+$ 并化简。
\end{proof}

\section{积分器与微分器}

\begin{proposition}[积分器]\label{prop:b1:operational-amplifier:integrator}
\index{积分器（运算放大器）}
输入经 $R$ 接至 $V_-$，以电容 $C$ 作反馈，$V_+$
接地：
\[
V_s(t) = -\frac{1}{RC}\int_0^t V_e(t')\,\dd t' + V_s(0),
\qquad
\underline H = -\frac{1}{jRC\omega} .
\]
在所有频率上，增益每十倍频程下降 \qty{20}{dB}，
\omterm{def:b1:wave-propagation:signal}{相位}为 $+90^\circ$：这是一个理想积分器——直到最小的失调或偏置
\omterm{def:b1:dc-circuits:voltage}{电压}被无限积分，最终把输出推入饱和。
\end{proposition}

\begin{proof}
\omterm{prop:b1:operational-amplifier:basic}{虚地}：电流 $V_e/R$ 流入电容，该电容的
\omterm{def:b1:dc-circuits:voltage}{电压}（从 $V_-$ 到输出端）为 $-V_s$；因此 $C\,\dd(-V_s)/\dd t
= V_e/R$，用复数表示为 $\underline H = -1/(jRC\omega)$。
\end{proof}

\begin{remark}[驯服积分器]\label{rem:b1:operational-amplifier:pseudo}
\index{伪积分器}
电阻 $R'$ 与 $C$ 并联，得
\[
\underline H = -\frac{R'/R}{1 + jR'C\omega},
\]
这是一阶低通\omterm{def:b1:filters-transfer-functions:H}{滤波器}，直流增益为 $-R'/R$，\omterm{def:b1:filters-transfer-functions:bode}{截止频率}为 $1/R'C$，它
在 $\omega \gg 1/R'C$ 时起积分作用，但输出不会漂移至饱和。类似地，\omterm{prop:b1:filters-transfer-functions:highpass1}{微分器}（输入端接电容，
反馈端接电阻：$V_s = -RC\,\dd V_e/\dd t$，$\underline H =
-jRC\omega$）会无限制地放大高频噪声，除非
在输入端串联一个小电阻以限制其增益。
\end{remark}

\begin{omfigure}
\begin{tikzpicture}[font=\small]
  \node[op amp] (c) at (0,0) {};
  \draw (c.+) -- ++(-0.4,0) node[ground]{};
  \draw (c.-) -- ++(-0.3,0) coordinate (m1) to[R=$R$] ++(-1.8,0) node[left] {$V_e$};
  \draw (c.out) -- ++(0.5,0) node[right] {$V_s$};
  \draw (c.out) ++(0.25,0) |- ++(0,1.2) to[C=$C$] ++(-2.1,0) -| (m1);
  \begin{scope}[xshift=7.3cm, yshift=-1.2cm]
  \begin{axis}[omaxis, axis x line=bottom, axis y line=left, xlabel style={at={(axis description cs:0.5,-0.16)}, anchor=north},
      width=6.4cm, height=4.2cm, xmode=log, xlabel={$\omega$ ($\log$)}, ylabel={$G_{\mathrm{dB}}$},
      xmin=0.01, xmax=100, ymin=-45, ymax=48, ytick={40,20,0,-20,-40}, xtick={0.01,0.1,1,10,100},
      xticklabels={$0.01$,$0.1$,$1$,$10$,$100$}, anchor=origin, at={(0,0)}]
    \addplot[omDef, thick, domain=0.01:100, samples=2] {-20*log10(x)};
    \addplot[omThm, thick, domain=0.01:100, samples=200] {20 - 10*log10(1 + (x/0.1)^2)};
    \node[omDef, font=\footnotesize, rotate=-32] at (axis cs:8,-10) {积分器};
    \node[omThm, font=\footnotesize, anchor=west] at (axis cs:0.012,29) {伪积分器};
  \end{axis}
  \end{scope}
\end{tikzpicture}

{\small 积分器及其\omterm{def:b1:filters-transfer-functions:bode}{伯德图}（蓝色，全程 \qty{-20}{dB/\text{十倍频程}}
直流处增益无穷大）；在 $C$ 两端\omterm{prop:b1:dc-circuits:seriesparallel}{并联电阻}后，低频增益
受到限制（红色）：这是一个在
\omterm{def:b1:filters-transfer-functions:bode}{截止频率}以上起积分作用的低通\omterm{def:b1:filters-transfer-functions:H}{滤波器}。}
\end{omfigure}

\begin{example}[由方波得三角波]\label{ex:b1:operational-amplifier:triangle}
$R = \qty{10}{k\ohm}$, $C = \qty{100}{nF}$, $1/RC = \qty{1000}{s^{-1}}$.
一个频率为 \qty{1}{kHz}、幅值为 $\pm\qty{1}{V}$ 的方波，得到斜率为
$\mp\qty{1000}{V/s}$、\omterm{def:b1:wave-propagation:signal}{振幅}为 $ET/4RC = \qty{0.25}{V}$ 的三角波——没有
无源 RC 电路所具有的 $1/\omega$ 衰减，且
输出\omterm{def:b1:sinusoidal-impedance:impedance}{阻抗}为零，可驱动下一级。
\end{example}

\section{有源滤波器}

\begin{proposition}[一阶有源低通滤波器]\label{prop:b1:operational-amplifier:activelp}
\index{有源滤波器}
\omterm{prop:b1:operational-amplifier:basic}{反相放大器}中，电容 $C$ 并联在
反馈电阻 $R_2$ 两端，其\omterm{def:b1:filters-transfer-functions:H}{传递函数}为
\[
\underline H = -\frac{R_2/R_1}{1 + jR_2C\omega}:
\]
这是一个低通\omterm{def:b1:filters-transfer-functions:H}{滤波器}，直流增益为 $-R_2/R_1$（其绝对值可大于 $1$），\omterm{def:b1:filters-transfer-functions:bode}{截止
频率}为 $\omega_c = 1/R_2C$，输出\omterm{def:b1:sinusoidal-impedance:impedance}{阻抗}为零——多级串接时，
各级\omterm{def:b1:filters-transfer-functions:H}{传递函数}直接相乘。
\end{proposition}

\begin{proof}
\omterm{prop:b1:operational-amplifier:basic}{反相放大器}，以反馈\omterm{def:b1:sinusoidal-impedance:impedance}{阻抗} $R_2 \parallel 1/jC\omega
= R_2/(1 + jR_2C\omega)$ 代替 $R_2$。
\end{proof}

\section{无负反馈：比较器}

\begin{proposition}[比较器]\label{prop:b1:operational-amplifier:comparator}
\index{比较器}
无反馈时，理想放大器处于饱和状态：
$V_s = +V_{\mathrm{sat}}$（若 $V_+ > V_-$），$-V_{\mathrm{sat}}$（若
$V_+ < V_-$）。一输入端接参考\omterm{def:b1:dc-circuits:voltage}{电压} $V_{\mathrm{ref}}$，
另一输入端接\omterm{def:b1:wave-propagation:signal}{信号}，输出即可判断\omterm{def:b1:wave-propagation:signal}{信号}是否超过
参考值——这是一个一位转换器。
\end{proposition}

\begin{proof}
$\mu\epsilon$ 超过 $V_{\mathrm{sat}}$，对任意 $\abs\epsilon >
V_{\mathrm{sat}}/\mu \sim \qty{0.1}{mV}$ 均成立。
\end{proof}

\begin{proposition}[施密特触发器]\label{prop:b1:operational-amplifier:schmitt}
\index{施密特触发器}\index{滞回}
将输出的一部分 $\beta = R_1/(R_1 + R_2)$ 反馈至 $V_+$
（$R_1$ 接地，$R_2$ 接输出端，构成分压），将\omterm{def:b1:wave-propagation:signal}{信号}
$V_e$ 接至 $V_-$。输出由 $+V_{\mathrm{sat}}$ 跳变为
$-V_{\mathrm{sat}}$，当 $V_e$ 上升超过 $+\beta V_{\mathrm{sat}}$ 时，
而 $V_e$ 下降低于 $-\beta V_{\mathrm{sat}}$ 时，输出跳回：两个阈值，
形成一个\emph{\omterm{prop:b1:operational-amplifier:schmitt}{滞回}}，宽度为 $2\beta V_{\mathrm{sat}}$，
小于该宽度的噪声不会引起输出变化。
\end{proposition}

\begin{proof}
\omterm{prop:b1:operational-amplifier:feedback}{正反馈}：饱和。若 $V_s = +V_{\mathrm{sat}}$，则 $V_+ =
\beta V_{\mathrm{sat}}$，且 $\epsilon = \beta V_{\mathrm{sat}} - V_e$
保持为正，直到 $V_e$ 超过 $\beta V_{\mathrm{sat}}$；此时
输出翻转，$V_+$ 降至 $-\beta V_{\mathrm{sat}}$，新状态
保持到 $V_e$ 降至该值以下。
\end{proof}

\begin{omfigure}
\begin{tikzpicture}[font=\small]
  \node[op amp] (s) at (0,0) {};
  \draw (s.-) -- ++(-0.8,0) node[left] {$V_e$};
  \draw (s.out) -- ++(0.6,0) node[right] {$V_s$};
  \coordinate (j) at ($(s.+)+(-0.6,0)$);
  \coordinate (f) at ($(j)+(0,-1.0)$);
  \coordinate (o) at ($(s.out)+(0.3,0)$);
  \draw (s.+) -- (j);
  \draw (o) -- (o |- f) to[R=$R_2$] (f) -- (j);
  \draw (f) to[R=$R_1$] ++(0,-1.6) node[ground]{};
  \begin{scope}[xshift=7.8cm, yshift=-0.9cm]
  \begin{axis}[omaxis, axis x line=bottom, axis y line=left, width=6.4cm, height=4.6cm, xlabel={$V_e$}, ylabel={$V_s$},
      xlabel style={at={(axis description cs:0.5,-0.12)}, anchor=north},
      xmin=-2.2, xmax=2.2, ymin=-1.5, ymax=1.5, xtick={-1,1}, xticklabels={$-\beta V_{\mathrm{sat}}$,$\beta V_{\mathrm{sat}}$},
      ytick={-1,1}, yticklabels={$-V_{\mathrm{sat}}$,$+V_{\mathrm{sat}}$}, anchor=origin, at={(0,0)}]
    \draw[omThm, thick, -Stealth] (axis cs:-2.2,1) -- (axis cs:0.2,1);
    \draw[omThm, thick] (axis cs:0.2,1) -- (axis cs:1,1) -- (axis cs:1,-1);
    \draw[omThm, thick, -Stealth] (axis cs:2.2,-1) -- (axis cs:-0.2,-1);
    \draw[omThm, thick] (axis cs:-0.2,-1) -- (axis cs:-1,-1) -- (axis cs:-1,1);
  \end{axis}
  \end{scope}
\end{tikzpicture}

{\small \omterm{prop:b1:operational-amplifier:schmitt}{施密特触发器}：经 $R_1$、$R_2$ 的\omterm{prop:b1:operational-amplifier:feedback}{正反馈}
给出两个阈值 $\pm\beta V_{\mathrm{sat}}$。输出
在 $V_e$ 上升超过上阈值时由高变低，
在下降至下阈值以下时由低变高（箭头所示），形成一个\omterm{prop:b1:operational-amplifier:schmitt}{滞回}循环。}
\end{omfigure}

\begin{example}[弛豫振荡器]\label{ex:b1:operational-amplifier:astable}
\index{弛豫振荡器}
将\omterm{prop:b1:operational-amplifier:schmitt}{施密特触发器}的输出接回其自身的反相输入端，
经 $R$ 接入，并将 $C$ 从 $V_-$ 接至地。
电容器向 $\pm V_{\mathrm{sat}}$ 充电，每次达到
$\pm\beta V_{\mathrm{sat}}$ 时使触发器翻转：输出为方波，
电容器\omterm{def:b1:dc-circuits:voltage}{电压}为一串指数弧段，周期为
\[
T = 2RC\ln\frac{1 + \beta}{1 - \beta}
\]
（见\cref{exo:b1:operational-amplifier:12}）：$2.2\,RC$，对应 $\beta = \tfrac12$。
无输入、一个电容器和一个时钟——一个系统，
因无法稳定而振荡，这一主题将在
本科二年级分册的反馈振荡器中再次出现。
\end{example}

\section{练习}

\begin{exercise}[$\star$]\label{exo:b1:operational-amplifier:1}
\omterm{prop:b1:operational-amplifier:basic}{同相放大器}中 $R_1 = \qty{1.0}{k\ohm}$，$R_2 =
\qty{9.0}{k\ohm}$，$V_{\mathrm{sat}} = \qty{15}{V}$。求增益？
当 $V_e = \qty{0.50}{V}$ 时的输出？饱和前的最大输入为多少？
\end{exercise}

\begin{exercise}[$\star$]\label{exo:b1:operational-amplifier:2}
\omterm{prop:b1:operational-amplifier:basic}{反相放大器}中 $R_1 = \qty{10}{k\ohm}$，$R_2 = \qty{100}{k\ohm}$。
求增益、输入\omterm{def:b1:sinusoidal-impedance:impedance}{阻抗}、从 \qty{1.0}{V} 电源汲取的电流以及输出。
\end{exercise}

\begin{exercise}[$\star$]\label{exo:b1:operational-amplifier:3}
一个\omterm{def:b1:dc-circuits:sources}{内阻}为 \qty{10}{k\ohm}、开路\omterm{def:b1:dc-circuits:voltage}{电压}为
\qty{1.0}{V} 的传感器需驱动一个 \qty{1.0}{k\ohm} 的负载。
在两者之间不接和接入跟随器两种情况下，负载两端的\omterm{def:b1:dc-circuits:voltage}{电压}各为多少？
\end{exercise}

\begin{exercise}[$\star$]\label{exo:b1:operational-amplifier:4}
\omterm{prop:b1:operational-amplifier:sumdiff}{求和放大器}中 $R_f = \qty{10}{k\ohm}$，两输入端分别经
\qty{10}{k\ohm} 和 \qty{20}{k\ohm} 接入。写出 $V_s$；并计算
$V_1 = \qty{1.0}{V}$、$V_2 = \qty{2.0}{V}$ 时的值。
\end{exercise}

\begin{exercise}[$\star\star$]\label{exo:b1:operational-amplifier:5}
推导\omterm{prop:b1:operational-amplifier:sumdiff}{差分放大器}的 $V_s = (R_2/R_1)(V_2 - V_1)$。取
$R_1 = \qty{10}{k\ohm}$，$R_2 = \qty{100}{k\ohm}$，$V_1 = \qty{1.00}{V}$，
$V_2 = \qty{1.05}{V}$：求输出？两输入端共有的
\qty{1}{V} \omterm{def:b1:dc-circuits:voltage}{电压}发生了什么变化？
\end{exercise}

\begin{exercise}[$\star\star$]\label{exo:b1:operational-amplifier:6}
积分器中 $R = \qty{10}{k\ohm}$，$C = \qty{100}{nF}$，输出
初始为零。施加一个 \qty{1.0}{V} 阶跃：求输出的斜率，
以及达到饱和（\qty{15}{V}）所需的时间。施加一个 \qty{1.0}{kHz} 的方波，
幅值为 $\pm\qty{1.0}{V}$：求输出的形状与\omterm{def:b1:wave-propagation:signal}{振幅}。
\end{exercise}

\begin{exercise}[$\star\star$]\label{exo:b1:operational-amplifier:7}
在前述积分器的电容器两端并联一个 \qty{1.0}{M\ohm} 的电阻。
求\omterm{def:b1:filters-transfer-functions:H}{传递函数}、直流增益、\omterm{def:b1:filters-transfer-functions:bode}{截止频率}，
以及 \qty{1.0}{kHz} 处的增益；
该电路在哪一频率范围内仍能积分？
\end{exercise}

\begin{exercise}[$\star\star$]\label{exo:b1:operational-amplifier:8}
设计一个反相有源低通\omterm{def:b1:filters-transfer-functions:H}{滤波器}，使其\omterm{def:b1:filters-transfer-functions:bode}{通带}增益为 \qty{20}{dB}，
\omterm{def:b1:filters-transfer-functions:bode}{截止频率}为 \qty{500}{Hz}，并取 $R_1 = \qty{1.0}{k\ohm}$。
\end{exercise}

\begin{exercise}[$\star\star$]\label{exo:b1:operational-amplifier:9}
一个实际放大器的 $\mu_0 = 10^5$，$f_T = \qty{1.0}{MHz}$。对于设计增益为
$100$ 的\omterm{prop:b1:operational-amplifier:basic}{同相放大器}：求实际直流增益
（分析中保留有限的 $\mu_0$），以及带宽；
跟随器的带宽是多少？
\end{exercise}

\begin{exercise}[$\star\star\star$]\label{exo:b1:operational-amplifier:10}
\omterm{prop:b1:operational-amplifier:comparator}{比较器}中 $V_{\mathrm{ref}} = \qty{2.0}{V}$ 接在 $V_-$ 端，用于监视
一个缓慢上升且含有 \qty{50}{mV} 噪声的\omterm{def:b1:wave-propagation:signal}{信号}。
描述\omterm{def:b1:wave-propagation:signal}{信号}越过 \qty{2.0}{V} 时输出的变化，并提出补救办法。
\end{exercise}

\begin{exercise}[$\star\star\star$]\label{exo:b1:operational-amplifier:11}
\omterm{prop:b1:operational-amplifier:schmitt}{施密特触发器}，$R_1 = \qty{1.0}{k\ohm}$，$R_2 = \qty{9.0}{k\ohm}$，
$V_{\mathrm{sat}} = \qty{15}{V}$，推导并计算两个阈值。
画出 $V_s$ 草图，其中 $V_e$ 为幅值 \qty{3}{V} 的三角波。
\end{exercise}

\begin{exercise}[$\star\star\star$]\label{exo:b1:operational-amplifier:12}
弛张振荡器，$R_1 = R_2$（$\beta = \tfrac12$）。推导
周期 $T = 2RC\ln\frac{1 + \beta}{1 - \beta}$，依据是 $C$ 经 $R$ 在两阈值之间的充电过程，
并选取 $R$，使频率为 \qty{1.0}{kHz}；
取 $C = \qty{100}{nF}$，画出 $V_s(t)$ 和 $V_C(t)$ 的草图。
\end{exercise}

\section{问题：应变片的测量链}

\begin{problem}[{周末问题——从桥梁主梁上的两毫伏信号
到一盏红灯：跟随器、差分放大器、有源滤波器
与施密特触发器，以及决定链路能否工作的
两个非理想因素}]
\label{pb:b1:operational-amplifier:1}
一个电阻为 $R(1 + \epsilon)$ 的应变片，$R = \qty{1.00}{k\ohm}$，
$0 \leq \epsilon \leq \num{2.0e-3}$，与三个 \qty{1.00}{k\ohm} 的固定电阻
组成\omterm{prop:b1:dc-circuits:wheatstone}{惠斯通电桥}，由 $E = \qty{5.00}{V}$ 供电。
电桥输出 $u = V_B - V_A$（两中点之间）为
$u \approx E\epsilon/4$（符号由接线决定）。放大器
除非另有说明均视为理想；$V_{\mathrm{sat}} = \qty{15}{V}$；必要时
取 $f_T = \qty{1.0}{MHz}$，\omterm{rem:b1:operational-amplifier:real}{压摆率}为 \qty{0.5}{V/\micro s}，\omterm{rem:b1:operational-amplifier:real}{输入失调电压}为
\qty{1}{mV}。

\medskip
\noindent\textbf{第一部分——电桥及其负载。}
\begin{enumerate}
  \item 回顾 $u \approx E\epsilon/4$ 在 $\epsilon$ 较小时成立的原因。
  \item 求 $u$ 在满量程（$\epsilon = \num{2.0e-3}$）时的值。
  \item 给出零应变与满量程时各中点的电势
        （相对于电桥的负接线柱），
        它们的共同值叫作什么？
  \item 每个中点都是由两个 \qty{1}{k\ohm} 电阻组成的\omterm{prop:b1:dc-circuits:dividers}{分压器}的输出，
        则两中点之间所见的戴维南电阻
        为多少？
  \item 若一台输入电阻为 \qty{10}{k\ohm} 的放大器
        接在电桥输出两端，
        则它接收到的是 $u$ 的几分之几？
  \item 为何在每个中点接一个跟随器就能解决这一问题？说明所用的
        两条性质。
  \item 跟随器的输出\omterm{def:b1:dc-circuits:voltage}{电压}（以 $V_A$、
        $V_B$ 表示）各为多少，它们从电桥汲取多大电流？
\end{enumerate}

\medskip
\noindent\textbf{第二部分——放大差分\omterm{def:b1:wave-propagation:signal}{信号}。}
\begin{enumerate}[resume]
  \item 跟随器的输出接入\omterm{prop:b1:operational-amplifier:sumdiff}{差分放大器}（两个输入端各接 $R_1$，
        $R_2$ 作反馈并接地）。推导 $V_s =
        (R_2/R_1)(V_B - V_A)$。
  \item 取 $R_1 = \qty{10}{k\ohm}$，并选取 $R_2$ 使增益为 $100$。
  \item 第二级为\omterm{prop:b1:operational-amplifier:basic}{同相放大器}，放大倍数为 $10$：给出其电阻，
        并求总增益及满量程时的输出。
  \item 两个中点都接近 $E/2 = \qty{2.5}{V}$（即共模\omterm{def:b1:dc-circuits:voltage}{电压}）。
        完全匹配的\omterm{prop:b1:operational-amplifier:sumdiff}{差分放大器}对共模\omterm{def:b1:dc-circuits:voltage}{电压}有何作用？
        若某一电阻比值偏差 $1\%$，估算折合到输入端的
        虚假\omterm{def:b1:dc-circuits:voltage}{电压}，并与\omterm{def:b1:wave-propagation:signal}{信号}相比较。
  \item 第一级放大器具有 \qty{1}{mV} 的\omterm{rem:b1:operational-amplifier:real}{输入失调电压}，
        经总增益放大后将产生多大的输出误差？并给出结论。
  \item 取 $f_T = \qty{1}{MHz}$，每一级的带宽是多少？
        对于低于 \qty{5}{Hz} 的\omterm{def:b1:wave-propagation:signal}{信号}，该带宽是否足够？
  \item 卡车驶近时，输出须在约 \qty{10}{ms} 内变化 \qty{2.5}{V}，
        \omterm{rem:b1:operational-amplifier:real}{压摆率}是否会构成限制？
\end{enumerate}

\medskip
\noindent\textbf{第三部分——滤波。}
放大后的\omterm{def:b1:wave-propagation:signal}{信号}中混有一个 \qty{50}{Hz} 的交流声，
它是从电网拾取的。
\begin{enumerate}[resume]
  \item 设计一个增益为 $-1$ 的反相有源低通\omterm{def:b1:filters-transfer-functions:H}{滤波器}，
        \omterm{def:b1:filters-transfer-functions:bode}{截止频率}为 \qty{10}{Hz}，取 $R_1 = R_2 = \qty{16}{k\ohm}$，求 $C$。
  \item 计算其在 \qty{50}{Hz} 处以\omterm{def:b1:filters-transfer-functions:bode}{分贝}表示的衰减量。
  \item 计算 \qty{5}{Hz} 处的相移及相应的
        延迟。对于报警器是否可接受？
  \item 此处为何宜选\omterm{prop:b1:operational-amplifier:activelp}{有源滤波器}而非无源 RC
        \omterm{def:b1:filters-transfer-functions:H}{滤波器}（两条理由）？
  \item 再级联一个相同的级，在 \qty{50}{Hz} 处的衰减量是多少？
        为何此处可以精确应用乘积法则？
\end{enumerate}

\medskip
\noindent\textbf{第四部分——阈值与报警。}
应变超过 \num{1.6e-3} 时，报警器必须动作。
\begin{enumerate}[resume]
  \item 这对应于测量链的哪一输出\omterm{def:b1:dc-circuits:voltage}{电压}？取该参考\omterm{def:b1:dc-circuits:voltage}{电压}的\omterm{prop:b1:operational-amplifier:comparator}{比较器}，
        其输出状态如何？
  \item 滤波后的\omterm{def:b1:wave-propagation:signal}{信号}在阈值附近仍含有 \qty{50}{mV} 的噪声，
        普通\omterm{prop:b1:operational-amplifier:comparator}{比较器}将如何动作？
  \item 一个\omterm{prop:b1:operational-amplifier:schmitt}{施密特触发器}，参考\omterm{def:b1:dc-circuits:voltage}{电压}接在一个输入端，反馈系数为
        $\beta$，证明其阈值为 $V_{\mathrm{ref}} \pm
        \beta V_{\mathrm{sat}}$；取 $\beta$ 使阈值为 $\pm\qty{0.10}{V}$，
        并给出 $R_1$、$R_2$。
  \item 分别以伏特和应变给出两个阈值。
  \item 输出经一只电阻驱动发光二极管（\qty{2.0}{V}，\qty{13}{mA}），
        求该电阻的阻值。输出处于
        $-V_{\mathrm{sat}}$ 时，发光二极管由什么来保护？
  \item 逐级概括整个测量链，并指出决定其实际极限的两个
        非理想因素。
\end{enumerate}
\end{problem}
